% GENERATED FILE. Edit canonical lessons, then run npm run generate:books.
% canonical-source-hash: fnv1a64:3e8360681ec46d10
% canonical-lessons: IT-C416-motorino, IT-C416-scooter, IT-C416-pedalare, IT-C416-parchimetro, IT-C416-sostare

\chapter{Su due ruote}
\label{ch:it-su-due-ruote}
\hlchaptermodality{416}

\begin{chapteropening}
\textbf{By the end of this chapter:} \emph{I can talk about getting round on two wheels and parking: the moped, the scooter, pedalling, the parking meter and stopping.}

Complete the last lesson of chapter 416: I can talk about getting round on two wheels and parking: the moped, the scooter, pedalling, the parking meter and stopping.
\end{chapteropening}

\section[il motorino]{il motorino — the moped}

\label{lesson:IT-C416-motorino}



\begin{quote}
Before the new one: say the Italian for the logbook, then the Italian for the insurance.
\end{quote}

\subsection*{You'll want to know: il motorino}
\textbf{il motorino} — \textquotedblleft{}the moped\textquotedblright{}.

\begin{quote}
Vado a scuola in motorino.
\end{quote}

\begin{quote}
\emph{I go to school by moped.}
\end{quote}

A small \emph{motore}; in many towns it is a teenager's first vehicle.

\begin{cousinweb}
\emph{motorino} is a little \emph{motore} 'engine', a learned borrowing from Latin \emph{mōtor} 'mover', from \emph{movēre} 'to move' (Lewis \& Short).
\end{cousinweb}

\begin{grammarlens}[title={What you've built}]
One more word for telling, explaining and describing at B1.
\end{grammarlens}

\subsection*{Your turn}
\begin{itemize}
\raggedright
  \item \emph{Say it:} \emph{il motorino}
  \item \emph{Say it:} \emph{il motorino}, once more
  \item \emph{Say it:} a sentence of your own with \emph{il motorino}
  \item \emph{Recall:} say the Italian for the logbook, then the Italian for the insurance, then say \emph{il motorino} again
\end{itemize}

\begin{tcolorbox}[breakable,colback=teal!4,colframe=teal!35!black,title={Before you move on}]
What does il motorino mean? (\textquotedblleft{}The moped\textquotedblright{}.) Say it once more.
\end{tcolorbox}
\section[lo scooter]{lo scooter — the scooter (motor scooter)}

\label{lesson:IT-C416-scooter}



\begin{quote}
Before the new one: say the Italian for the insurance, then the Italian for the moped.
\end{quote}

\subsection*{You'll want to know: lo scooter}
\textbf{lo scooter} — \textquotedblleft{}the scooter (motor scooter)\textquotedblright{}.

\begin{quote}
Con lo scooter giro la città in dieci minuti.
\end{quote}

\begin{quote}
\emph{With the scooter I get round the city in ten minutes.}
\end{quote}

Bigger than \emph{il motorino}; plural \emph{gli scooter}.

\begin{cousinweb}
\emph{scooter} was borrowed from English \emph{scooter} (EtymDB), built on \emph{scoot} 'to dart off'.
\end{cousinweb}

\begin{grammarlens}[title={What you've built}]
One more word for telling, explaining and describing at B1.
\end{grammarlens}

\subsection*{Your turn}
\begin{itemize}
\raggedright
  \item \emph{Say it:} \emph{lo scooter}
  \item \emph{Say it:} \emph{lo scooter}, once more
  \item \emph{Say it:} a sentence of your own with \emph{lo scooter}
  \item \emph{Recall:} say the Italian for the insurance, then the Italian for the moped, then say \emph{lo scooter} again
\end{itemize}

\begin{tcolorbox}[breakable,colback=teal!4,colframe=teal!35!black,title={Before you move on}]
What does lo scooter mean? (\textquotedblleft{}The scooter (motor scooter)\textquotedblright{}.) Say it once more.
\end{tcolorbox}
\section[pedalare]{pedalare — to pedal, to cycle}

\label{lesson:IT-C416-pedalare}



\begin{quote}
Before the new one: say the Italian for the moped, then the Italian for the scooter.
\end{quote}

\subsection*{You'll want to know: pedalare}
\textbf{pedalare} — \textquotedblleft{}to pedal, to cycle\textquotedblright{}.

\begin{quote}
Abbiamo pedalato per venti chilometri.
\end{quote}

\begin{quote}
\emph{We cycled twenty kilometres.}
\end{quote}

In speech, \emph{pedala!} also means get moving!

\begin{cousinweb}
\emph{pedalare} is built on \emph{pedale} 'pedal', a learned borrowing from Latin \emph{pedālis} 'of the foot' (EtymDB), from \emph{pēs} 'foot'.
\end{cousinweb}

\begin{grammarlens}[title={What you've built}]
One more word for telling, explaining and describing at B1.
\end{grammarlens}

\subsection*{Your turn}
\begin{itemize}
\raggedright
  \item \emph{Say it:} \emph{pedalare}
  \item \emph{Say it:} \emph{pedalare}, once more
  \item \emph{Say it:} a sentence of your own with \emph{pedalare}
  \item \emph{Recall:} say the Italian for the moped, then the Italian for the scooter, then say \emph{pedalare} again
\end{itemize}

\begin{tcolorbox}[breakable,colback=teal!4,colframe=teal!35!black,title={Before you move on}]
What does pedalare mean? (\textquotedblleft{}To pedal, to cycle\textquotedblright{}.) Say it once more.
\end{tcolorbox}
\section[il parchimetro]{il parchimetro — the parking meter}

\label{lesson:IT-C416-parchimetro}



\begin{quote}
Before the new one: say the Italian for the scooter, then the Italian for to pedal.
\end{quote}

\subsection*{You'll want to know: il parchimetro}
\textbf{il parchimetro} — \textquotedblleft{}the parking meter\textquotedblright{}.

\begin{quote}
Il parchimetro accetta solo monete.
\end{quote}

\begin{quote}
\emph{The parking meter takes coins only.}
\end{quote}

You pay and leave the ticket behind the windscreen: \emph{dietro il parabrezza}.

\begin{cousinweb}
\emph{parchimetro} is built from \emph{parcheggio} 'parking' and \emph{-metro}, a learned element that goes back through Latin \emph{metrum} to Greek \emph{métron} 'measure' (EtymDB): it measures your parking time.
\end{cousinweb}

\begin{grammarlens}[title={What you've built}]
One more word for telling, explaining and describing at B1.
\end{grammarlens}

\subsection*{Your turn}
\begin{itemize}
\raggedright
  \item \emph{Say it:} \emph{il parchimetro}
  \item \emph{Say it:} \emph{il parchimetro}, once more
  \item \emph{Say it:} a sentence of your own with \emph{il parchimetro}
  \item \emph{Recall:} say the Italian for the scooter, then the Italian for to pedal, then say \emph{il parchimetro} again
\end{itemize}

\begin{tcolorbox}[breakable,colback=teal!4,colframe=teal!35!black,title={Before you move on}]
What does il parchimetro mean? (\textquotedblleft{}The parking meter\textquotedblright{}.) Say it once more.
\end{tcolorbox}
\section[sostare]{sostare — to stop briefly, to park (for a short time)}

\label{lesson:IT-C416-sostare}



\begin{quote}
Before the new one: say the Italian for to pedal, then the Italian for the parking meter.
\end{quote}

\subsection*{You'll want to know: sostare}
\textbf{sostare} — \textquotedblleft{}to stop briefly, to park (for a short time)\textquotedblright{}.

\begin{quote}
Qui è vietato sostare.
\end{quote}

\begin{quote}
\emph{No stopping here.}
\end{quote}

\emph{Divieto di sosta} — no parking; \emph{una sosta} is a short stop.

\begin{cousinweb}
\emph{sostare} is inherited from Latin \emph{substāre} 'to stand firm, to hold out', built from \emph{sub} and \emph{stāre} 'to stand' (Lewis \& Short).
\end{cousinweb}

\begin{grammarlens}[title={What you've built}]
One more word for telling, explaining and describing at B1.
\end{grammarlens}

\subsection*{Your turn}
\begin{itemize}
\raggedright
  \item \emph{Say it:} \emph{sostare}
  \item \emph{Say it:} \emph{sostare}, once more
  \item \emph{Say it:} a sentence of your own with \emph{sostare}
  \item \emph{Recall:} say the Italian for to pedal, then the Italian for the parking meter, then say \emph{sostare} again
\end{itemize}

\begin{tcolorbox}[breakable,colback=teal!4,colframe=teal!35!black,title={Before you move on}]
What does sostare mean? (\textquotedblleft{}To stop briefly, to park (for a short time)\textquotedblright{}.) Say it once more.
\end{tcolorbox}
